%%
%% Cover Letter — Agricultural Systems (Elsevier)
%%
\documentclass[12pt]{article}
\usepackage[T1]{fontenc}
\usepackage[utf8]{inputenc}
\usepackage[american]{babel}
\usepackage[margin=2.5cm]{geometry}
\usepackage{setspace}
\onehalfspacing
\usepackage{hyperref}

\pagestyle{empty}

\begin{document}

\noindent
Cover Letter\\
\textit{Agricultural Systems}\\
Elsevier

\vspace{18pt}

\noindent
Dear Editors,

\vspace{12pt}

\noindent
On behalf of all co-authors, I am pleased to submit the manuscript entitled ``Machine Learning Readiness to Support Agroecological Transitions in Traditional Agricultural Systems'' for consideration in \textit{Agricultural Systems}.

\vspace{6pt}

\noindent
This work presents the first comprehensive meta-analysis assessing the operational readiness of machine learning to support evidence-based monitoring and agroecological transitions in Traditional Agricultural Systems (TAS). Analyzing 244 studies (2010--2025) through random-effects modeling, bibliometric network analysis, Multiple Correspondence Analysis, and FAIR compliance assessment, we quantify the gap between reported algorithmic performance and actual deployment readiness for regulatory and heritage-management applications. Key findings include: (i)~pooled classification accuracy of 90.7\% (95\% CI 89.8--91.5\%) with substantial heterogeneity ($I^{2} = 58\%$); (ii)~deep learning adoption rising from 4.8\% to 21.1\% between 2015--2019 and 2020--2025; (iii)~critically low FAIR compliance (mean 18.7/100), with only 1.1\% of studies sharing code or data; and (iv)~publication-bias diagnostics (Egger's test, $p = 0.009$; \textit{trim-and-fill} correction of $-$1.54~pp) indicating systematic overestimation.

\vspace{6pt}

\noindent
We believe this manuscript aligns closely with the scope of \textit{Agricultural Systems}. The study adopts a systems-analysis perspective to evaluate how machine learning and remote sensing technologies interact with the biophysical, institutional, and data-governance dimensions of Traditional Agricultural Systems, examining 244 studies (2010--2025) through random-effects meta-analysis, co-occurrence network analysis, Multiple Correspondence Analysis, and FAIR compliance assessment. These systems, recognized by FAO's GIAHS program and by IPBES assessments, are coupled social-ecological systems whose management, monitoring, and policy evaluation depend on integrating heterogeneous data sources across spatial and temporal scales. Our analysis of spatial validation gaps, algorithmic transferability across farming contexts, and data-sharing practices directly addresses the journal's interest in understanding how analytical tools and information systems shape the performance and governance of agricultural systems at multiple scales. The FAIR compliance evaluation reveals that near-absent code and data sharing (1.1\%) and low interoperability compromise the reproducibility and regulatory applicability of machine learning pipelines, a finding with direct implications for decision-support infrastructure in agricultural systems research. By bridging systems-level performance assessment, computational intelligence, and data-governance analysis, the manuscript contributes to the interdisciplinary scope of \textit{Agricultural Systems} at the interface of farming systems research, quantitative methods, and technology adoption for sustainable agricultural management.

\vspace{6pt}

\noindent
The manuscript represents original, unpublished research not under consideration elsewhere. All authors have approved the final version and consent to submission. All reproducibility artifacts, including the bibliographic corpus, screening scripts, and analytical routines, have been publicly deposited on the Open Science Framework (OSF) under DOI \url{https://doi.org/10.17605/OSF.IO/J7STC}, ensuring full transparency and replicability.

\vspace{6pt}

\noindent
We thank you for your consideration and look forward to your feedback.

\vspace{36pt}

\noindent
Sincerely,

\vspace{36pt}

\noindent
Luiz Diego Vidal Santos, PhD\\
(Corresponding Author)\\
Universidade Estadual de Feira de Santana (UEFS)\\
E-mail: ldvsantos@uefs.br

\end{document}
